\documentclass[10pt]{article}

\usepackage[utf8]{inputenc}

\usepackage[T1]{fontenc}

\usepackage{amsmath}

\usepackage{amsfonts}

\usepackage{amssymb}

\usepackage[version=4]{mhchem}

\usepackage{stmaryrd}



\begin{document}

пам (эллиптическому, гиперболическому или параболическому).



Линейное (или квазилинейное) дифференциальное уравнение 2-го порядка с двумя неизвестными переменными



\[

A u_{x x}+2 B_{x y}+C u_{y y}=f\left(x, y, u, u_{x}, u_{y}\right)

\]



Iі с непрерывными коәффициентами в области задания $\Omega$ является С. т. у., если в этой области дискриминант $\Delta=A C-B^{2}$ характеристич. формы



\[

A d y^{2}+2 B d x d y+C d x^{2} \equiv Q

\]



обращается в нуль, не будучи там тождественно равным нулю.



Кривая $\delta$, определяемая уравнением $\Delta=0$, наз. п а р а б о л ч ч е с к о й л ин и е й уравнения (1), или ли ние й выро ждения (из мен е ния) ти п а у а внения.



Если дискриминант $\Delta$ в области $\Omega$ не меняет знака при переходе точки $(x, y)$ через параболич. линию $\delta$, то уравнение (1) относится к вырожденным уравнениям әл ли ит ик о - п а р а б о личес к о го $(\Delta \geqslant$ $\geqslant 0)$ или ги пе р боло-па р а болическ ого $(\Delta \leqslant 0)$ т и п а (см. Вырожденное уравнение с частными производными).



При нек-рых условиях гладкости коэффициентов $A, B, C$ и параболич. линии $\delta$ существуют неособое действительное преобразование независимых переменных, Іриводящее уравнение (1) со знакопеременным дискриминантом $\Delta$ (в окрестности выбранной точки линии $\delta$, где $A^{2}+B^{2}+C^{2} \neq 0$ ) к одному из следующих канонич. видов (обозначения для независимых переменных сохранены):



\[

\begin{aligned}

& y^{2 m+1} u_{x x}+u_{y y}=F\left(x, y, u, u_{x}, u_{y}\right), \\

& u_{x x}+y^{2 m+1} u_{y y}=F\left(x, y, u, u_{x}, u_{y}\right) .

\end{aligned}

\]



У равнения (2) и (3) являются С. т. У. (эллиптико-параболич. типа) в любой области, содержащей внутри себя интервал линии вырождения $y=0$.



Область $\Omega$ задания С. т. У. принято называть с м еш н н о й о бл а с т ю ю, а краевые задачи в смешанных областях - с м е ші а н н м и к р а е в ым и $а$ а д а ч а м и. Часть $\Omega^{+}\left(\Omega^{-}\right)$смешанной области $\Omega$, где уравнение принадлежит эллиптическому (гиперболическому) типу, наз. о б л а с т ь ю ә л л и птичности (ги пе р боличности).



K отысканию определенных решений С. Т. У. сводятся многие проблемы прикладного характера, в тастности проблемы околозвукового течения сжимаемой среды и безмоментной теории оболочек.



C. т. у. (1) наз. уравнением п е р в о г о р о д а (в т о о го ро д а), если всюду вдоль параболич. линии характеристич. форма $Q \neq 0 \quad(Q=0)$. У р а в н ени е Ч а плы гин а



\[

k(y) u_{x x}+u_{y y}=0,

\]



где $k(y)$ - непрерывно дифференцируемая монотонная функция такая, что $y k(y)>0$ при $y \neq 0$, - типичный пример С. т. У. 1-го рода. При $k(y)=y$ уравнение (4) принято называть трикоми уравнением.



Важной моделью С. т. у. (с разрывными коәффициентами при старших производных) является у р а внени е Ја в рентьев а- Бицад 3 е



\[

\operatorname{sign} y u_{x x}+u_{y y}=0 \text {. }

\]



Одной из основных краевых задач для С. т. у. (первого рода) является з а д а ч а Т р и ко м и, к-рая для уравнения вида (2) ставится следующим образом. Пусть $\Omega$ - конечная односвязная область евклидової плоскости независимых переменных $x$ и $y$, ограниченная простой жордановой кривой $\sigma$ с концами в точках $A(0,0), B(1,0)$, лежащей в полуплоскости $y>0$, и частями $A C$ и $B C$ характеристик уравнения (2), выходящими из точки $C\left(1 / 2, y_{C}\right), y_{C}<0$. Задача Трикоми заключается в отыскании решения $u(x, y)$ уравнения (2), непрерывного в замыкании $\bar{\Omega}$ области $\Omega$ и принимающего наперед заданные значения на кривої $\sigma U A C$.



В теории задачи Трикоми существенную роль играет п р и ци п әкст р м ум а Бицад к-рый в случае уравнения (5) гласит: репение $u(x, y)$ уравнения (5) из класса $C \overline{(\Omega)} \cap C^{1}(\Omega)$, обращающееся в нуль на характеристике $A C^{\prime}: x+y=0,0 \leqslant x \leqslant 1 / 2$, в замыкании $\bar{\Omega}+$ области эллиптичности $\Omega+=\Omega \cap\{y>0\}$ своего экстремума достигает на кривой $\sigma$.



Этот принцип, из к-рого следует единственность и устойчивость решения задачи Трикоми, а также обоснование альтернирующего метода его отыскания, распространен на весьма широкий класс линейных и квазилинейных С. т. у. В частности, этому классу принадлежат уравнения Чаплыгина (и Трикоми), если $k(y)$ дважды непрерывно дифференцируема и $5 k^{\prime 2} \geqslant 4 k k^{\prime}$ при $y<0$. Принцип әкстремума Бицадзе остается в силе и для уравнения



\[

\text { sign } y \cdot|y|^{\alpha} u_{x x}+u_{y y}=0, \alpha=\text { const }>0 \text {. }

\]



Решение задачи Трикоми для уравнения (6) в соответствующей смешанной области $\Omega$ выписывается в явном виде, если әллиптич. часть $\sigma$ границы этой области совшадает с т. н. н о р м а ь ны м конт у р о м $\sigma_{0}$ :



\[

x^{2}+\left(\frac{2}{\alpha+2}\right)^{2} y^{\alpha+2}=1 / 4 .

\]



В общем случае при определенных условиях на кривую $\sigma$ и на класс искомых решений задача Трикоми для уравнения (6) эквивалентно редуцируется к сингулярному интегральному уравнению, безусловная разрешимость к-рого следует из единственности решения. Метод интегральных уравнений успешно применяется и при доказательстве существования решения задачи Трикоми и др. смешанных задач для более общих уравнений вида



\[

\operatorname{sign} y \cdot|y|^{\alpha} u_{x x}+u_{y y}=F\left(x, y, u, u_{x}, u_{y}\right)

\]



со степенным вырождением порядка $\alpha$.



Методы теории функций и функционального анализа, особенно метод ашриорных оценок, позволили значительно расширить класс С. т. У. и смешанных областей, для к-рых имеет место единственность и существование (обобщенного) решения как задачи Трикоми, так и ряда др. смешанных задач.



Существенным обобщением задачи Трикоми является о бщая сме шш анная з а д а ч а Б иц а д з е, к-рая в случае уравнения (5) ставится следующим образом. Пусть $\Omega$ - односвязная смешанная область, ограниченная лежащей в полуплоскости $y>0$ простой жордановой кривой $\sigma$ с концами в точках $A(0$, $0), B(1,0)$ и выходящими из этих точек (гладкими) монотонными кривыми $\Gamma_{0}$ и $\Gamma_{1}$, к-рые пересекаются в точке $C\left(x_{1}, y_{1}\right), y_{1}<0$. Предполагается, что кривые $\Gamma_{0}$ и $\Gamma_{1}$ принадлежат области, ограниченной характеристиками $x+y=0, x-y=1$ и отрезком $A B: 0 \leqslant x \leqslant 1$ прямой $y=0$. Через $B_{0}$ и $B_{1}$ обозначены точки пересечения характеристик $x-y_{1}=x_{0}$ и $x+y=x_{0}$ с кривыми $\Gamma_{0}$ и $\Gamma_{1}$, где $x_{0}$ - любая фиксированная точка из полуинтервала $x_{1}+y_{1}<x_{1} \leqslant x_{1}-y_{1}$, а через $\gamma_{0}$ и $\gamma_{1}-$ части кривых $\Gamma_{0}$ и $\Gamma_{1}$, лежащих между точками $A, B_{0}$ и $B, B_{1}$ соответственно. Общая смешанная задача Бицадзе заключается в отыскании регулярного (при $y \neq 0$, $x \pm y \neq x_{0}$ ) решения уравнения (5) в области $\Omega$, к-рое непрерывно в $\bar{\Omega}$, имеет непрерывные первые производные в $\Omega$ гри $x \pm y \neq x_{0}$ и удовлетворяет заданным краевым условиям на кривых $\sigma, \gamma_{0}$ и $\gamma_{1}$. Единствен-





\end{document}